% Generated draft: requires mathematical and source review.
\documentclass[11pt]{article}
\usepackage[margin=25mm]{geometry}
\usepackage{fontspec}
\setmainfont{Latin Modern Roman}
\usepackage{amsmath,amssymb,mathtools,unicode-math}
\setmathfont{Latin Modern Math}
\usepackage{enumitem,needspace,xurl,hyperref}
\setlength{\parindent}{0pt}
\setlength{\parskip}{5pt}
\setlength{\emergencystretch}{3em}
\newcommand{\N}{\mathbb{N}}
\newcommand{\Z}{\mathbb{Z}}
\newcommand{\Q}{\mathbb{Q}}
\newcommand{\R}{\mathbb{R}}
\newcommand{\C}{\mathbb{C}}
\newcommand{\F}{\mathbb{F}}
\newcommand{\A}{\mathbb{A}}
\newcommand{\PP}{\mathbb{P}}
\newcommand{\E}{\mathbb{E}}
\newcommand{\eps}{\varepsilon}
\newcommand{\dd}{\,\mathrm{d}}
\begin{document}
\section*{UnsolvedMath ID: 7000004}
\textbf{Source problem code:} catalogue entry
\subsection{Definitions and mathematical statement}
Represent the smooth closed immersed curve by a regular map $\gamma:S^1\to\R^3$. At points of nonzero curvature its usual unit binormal is
\[
 B(s)=\frac{\gamma'(s)\times\gamma''(s)}
            {|\gamma'(s)\times\gamma''(s)|}.
\]
The source assumes a continuous binormal field along the whole curve and requires it to be one-to-one. Under the standard unit-field and parameter-circle reading this means $B:S^1\to S^2$ is injective. The record does not state a curvature-nonvanishing hypothesis, so any extension across zero-curvature points needs the source's convention rather than exclusion of those points.

The associated thin ribbon has the local form
\[
 F(s,u)=\gamma(s)+uB(s),\qquad |u|\text{ small},
\]
with the regularity and embedding convention to be fixed. The question is whether the injectivity of $B$ forces this ribbon $(\Gamma,B)$ to be ``twisted.'' The source does not define twisted. A total-twist condition, nontrivial framing, linking of boundary curves, and nonorientability are different notions; no one is imposed here. In particular the curve is only immersed, so a definition requiring disjoint embedded boundary curves cannot automatically be used.

The exact retained target is therefore the implication from a continuous injective binormal field to the source's twisted-ribbon property, pending that definition. Unit normalization, whether injectivity is on the parameter circle or geometric image, and the meaning of binormal at inflections must also be confirmed. The spherical image specified by the actual hypothesis is the binormal image, not a silently substituted principal-normal indicatrix.
\subsection{Short English statement}
Does a smooth closed immersed space curve with a continuous one-to-one binormal field necessarily determine a twisted ribbon? The central missing input is the intended notion of ribbon twist, especially since the curve need not be embedded.

The supplied summary mentions special curve classes, thin strips, and width/twist inequalities but gives no exact theorem or direct solved reference. It calls the general assertion unresolved. Its references to a normal spherical image do not replace the stated binormal condition. The partial-progress label and 2019 survey problem are retained with current status marked for review.
\subsection{Sources}
\begin{itemize}
\item Local UnsolvedMath registry, permanent ID 7000004, source code catalogue entry, ``Geometry of Curves and Surfaces --- Problem 1.4,'' proposed by Mohammad Ghomi, 2019. Dataset curation and metadata: Ulam AI and UnsolvedMath Contributors, CC BY 4.0.
\item Mohammad Ghomi, \emph{Open Problems in Geometry of Curves and Surfaces} (2019), Problem 1.4, PDF page 6, \url{https://people.math.gatech.edu/~ghomi/Papers/op.pdf}; normalized pdf-text extraction, accessed 2026-07-29.
\item The supplied summary mentions Ghomi's width work and special curve classes without separate full references or exact twist results. Source-specific provenance: the author survey labels the item a problem; current status and statement-reuse terms NEEDS\_REVIEW for the public author-hosted PDF, separately from the dataset licence.
\end{itemize}
\end{document}
